\section{表面}
理想周期体系是永恒的。但实际上，体相体系的极限是其表面。\textrm{VASP}模拟表面时，在两个维度上保持体相结构，在一个方向上引入真空，从而构成表面\textrm{(surface)}或厚板\textrm{(slab)}。学习表面的结构弛豫，计算\textrm{DOS}和能带，模拟表面的分子吸附和\textrm{STM}图像。

\subsection{\ch{Ni}表面}
对两个五层\textrm{(100)}\ch{Ni}表面执行自旋极化计算，采用\textrm{RMM-DIIS}方法完成结构弛豫。通过对比表面能和体相计算得到弛豫能。

表面是实际体相体系的最前沿，\textrm{VASP}模拟表面时，将晶胞沿某个方向延伸结构，使得原有材料部分保持体相结构，部分则是被真空填充。这样的结构称为表面或\textrm{Slab}或直接称呼\textrm{Slab}。一般沿延伸结构方向用一个$\vec k$点表示，表明真空区域隔绝了重复单元间的相互作用。增加而外的$\vec k$点，有助于描述相邻晶胞的相互作用。\textcolor{red}{这样的重复单元足够厚(包含真空区和材料晶胞区域同样要求)，才可以在\textrm{Slab}中间，继续保持体相性质的模拟。换言之，沿着厚度延伸方向，不能再依赖周期边界条件；而模拟表面时，必须消除真空延伸带来的周期性影响。}

\ch{Ni}-\textrm{(100)}面的\textrm{POSCAR}文件

体相\textrm{POSCAR}

\textrm{INCAR}

体相\textrm{INCAR}

$\vec k$空间剖分网格方案\textrm{KPOINTS}

体相\textrm{KPOINTS}

\textrm{POSCAR}中确定的晶格常数为\textrm{3.53\AA}。体相晶胞中只有一个原子，在表面计算中，每层一个原子，共计有五层。真空层大小为$3\times3.53\AA=10.59\AA$。采用\texttt{selective dynamics}模式允许两层原子弛豫而剩下三层原子固定。

根据自旋极化计算的要求，设置\texttt{ISPIN}=2且有\texttt{MAGMOM}=1。%换言之，在\textrm{Kohn-Sham}轨道求解的时候，两部分自旋电子是分别考虑的，而构造\textrm{Hamiltonian}时，则要系统考虑全部电子。
初始电荷密度通过原子电荷密度叠加得到。%\texttt{ICHARG}设置使得占据波函数由\textrm{Methfessel-Paxton}方法近似得到\upcite{PRB40-3616_1989}。
结构弛豫由计算控制参数\texttt{IBRION}，\texttt{NSW}、\texttt{POTIM}和\texttt{ISIF}确定。为了使得原子受到的力和应力为零，在给定赝势基础上，依次执行
\begin{itemize}
	\item 给定原子位置，通过\textrm{SDFT}计算电子基态；
	\item 利用\textrm{Hellmann-Feyman}定理，计算电子\textrm{Hamiltonian}的梯度值，作为作用于原子上的力和应力；
	\item 采用\textrm{RMM-DIIS}方法，离子弛豫到其瞬时基态；
	\item 重复上述三步，直到体系达到收敛标准。
\end{itemize}
在\textrm{KPOINTS}中，$\vec k$空间网格点由\textrm{Monkhorst-Pack}方法确定。$\vec k$空间网格布点确保$\Gamma$点在中心位置。对于表面计算，在表面方向，只有一个$\vec k$点。

分别对体相和表面执行\textrm{Bash}脚本，完成计算。

采用\textrm{Py4vasp}模块可视化体相结构，构建$3\times3\times3$的晶胞。

相比于体相的结构弛豫，表面的结构弛豫在真空方向(本算例中的$z$方向)

表面的弛豫

采用\textrm{Py4vasp}模块检查表面计算的收敛情况。

\textrm{OUTCAR}文件中确认每个离子步的力的计算情况。第一个和最后一个离子步结束，原子受力有何不同，是否达到\texttt{EDIFFG}设置的收敛标注。

在第一个离子步，第一、第二层原子的受力与第四、第五层原子是相当的。\textrm{POSCAR}中的控制参数\texttt{selective dynamics}许可第四、第五层原子弛豫。

到第五个离子步，第四、第五层原子的受力已经极小。$\mathtt{EDIFFG}=10\times\mathtt{EDIFF}=1E-5$。

\textcolor{red}{最后计算表面能}
非弛豫表面能
\begin{displaymath}
	\begin{aligned}
		\sigma^{\mathrm{unrel}}=&\dfrac1{2A}(E_{\mathrm{slab}^{\mathrm{unrel}}}-NE_{\mathrm{bluk}})\\
		=&\dfrac1{2\times3.53^2/2}[-25.5567-5\cdot(-5.4681)]\\
		=&143.0\mathrm{meV/\AA^2}
	\end{aligned}
\end{displaymath}
这里利用了非弛豫的\textrm{Slab}的总能$E_{\mathrm{slab}}^{\mathrm{unrel}}=-25.5567\mathrm{eV}$。

表面能的平均值则是通过计算弛豫\textrm{Slab}一侧上两层(第四、第五层)的能量得到：
\begin{displaymath}
	\sigma^{\mathrm{avg}}=\dfrac1{2A}(E_{\mathrm{slab}}-NE_{\mathrm{bulk}})=141.5\mathrm{meV/\AA^2}
\end{displaymath}

最后得到弛豫的表面能为：
\begin{displaymath}
	\sigma^{\mathrm{rel}}=2\sigma^{\mathrm{avg}}-\simgma^{\mathrm{unrel}}=140.1\mathrm{eV/\AA^2}
\end{displaymath}
可见表面弛豫能非常小。差不多是$3\mathrm{meV/\AA^2}$。但有些材料模拟中，弛豫能计算方法相同，不过弛豫的原子数比\textrm{FCC}密堆积金属\ch{Ni}的要多。文献给出100个\ch{Ni}原子的\textrm{PBE}表面能是$138\mathrm{meV/\AA^2}$\upcite{SD3-16080_2016}。可见当前的部分弛豫方法使用了最小的计算开销，但也能得到比较好的结果。

主要到表面能的计算表达式
\begin{displaymath}
	\sigma=\dfrac{E_{\mathrm{slab}}-NE_{\mathrm{bluk}}}{2A}
\end{displaymath}
只对具有一定对称性和正确化学计量比的表面有效。公式中有因子$1/2$是因为当引入真空体相形成表面的时候，就同时存在两个表面。这个公式成立的条件就是两个表面必须是等价的，也就意味着\textrm{slab}是对称的。系数$N$表明\textrm{Slab}包含$N$个构成体相的最小晶胞数目，所以只有当\textrm{Slab}中的元素的化学计量比于体相元素时完全等同时才成立。

由于体相材料中原子(或离子)因为完全成键，化学环境稳定。而表面原子与体相原子化学环境不同，一般都含有悬挂键(通常是未成对的自由基)，因此具有的能量高于体相原子。对照实验结果，\textrm{DFT}计算中，采用半局域交换-相关泛函计算表面体系，往往会高估其稳定性。也就是说，理论模拟得到的表面能，往往比实验值要低。

另外，根据表面结构的一般规律，$(211)$面的表面能比$(111)$面的表面能要高，这表明构成$(211)$面需要更多的能量，换言之，对于\textrm{FCC}密堆积体系，沿$(211)$面切出表面比$(111)$面更困难。

\subsubsection{\ch{Ni}的(100)表面的局域态密度}
对\ch{Ni}的$(100)$面，计算每个\ch{Ni}原子的电荷和局域\textrm{DOS}。

为了计算每个\ch{Ni}原子的电荷，必须首先确认哪些电荷属于哪些原子。%或者更一般地，需要通过哪些局域\textrm{Kohn-Sham}轨道计算哪些局域电荷。

一个计算方案是划定一个小的范围，规定在此范围内的电子属于该原子，这也就是\textrm{POTCAR}文件中给出参数\texttt{RWIGS}，定义了元素的\textrm{Wigner-Seitz}半径。当计算控制参数\texttt{LORBIT}$<10$，可以根据\textrm{Wigner-Seitz}半径，计算原子周围的电荷密度。该方案的问题是电荷密度的划分较为笼统，没有办法进一步根据角动量$l$、磁量子数$m$和自旋量子数$\sigma$来进一步区分电荷密度。

另一个方案是对电荷密度执行幺正变换，也就是对离域的\textrm{Kohn-Sham}轨道的子空间(占据轨道)实施变换，形成一套局域态基组，这样得到的基函数称为\textrm{Wannier}函数。显然\textrm{Wannier}函数并不唯一，通常说的\textrm{Wannier}函数是最大局域\textrm{Wannier}函数。如果\textrm{Kohn-Sham}轨道是强相干的\textrm{(strongly entangled)}，实施这种幺正变换会非常困难。

还有一种办法，是采用投影子算符\textrm{(Projector operators)}，将\textrm{Kohn-Sham}轨道的某些特征部分投影出来。这样构成的基函数就是投影局域轨道\textrm{(Projected localized orbitals)}\textcolor{red}{\upcite{PRB77-205112_2008}}。这就是\textrm{INCAR}文件中计算控制参数\texttt{LORBIT}$>10$时，计算投影出来的局域分波\textrm{DOS}。投影用的投影子和相因子系数也都写入\textrm{PROCAR}文件。

\ch{Ni}的$(100)$表面的结构文件\textrm{POSCAR}

\textrm{INCAR}

$\vec k$空间剖分网格文件\textrm{KPOINTS}

其中\textrm{POSCAR}文件来自前面计算的结构弛豫文件\textrm{CONTCAR}。

\textrm{INCAR}文件中，计算控制参数\texttt{ISMEAR}$=-5$，\texttt{LORBIT}$=11$。%\texttt{ISMEAR}=-5将使用\textrm{Bl\"ochl}校正的四面体积分方法，该方法对于分数占据的轨道是不可变分的，因此对于存在部分占据轨道(金属或导体)，计算力和应力会有误差。所以不能用于导体和金属的几何结构优化。\texttt{LORBIT}=11，关于\textrm{DOS}的计算结果将同时写入\textrm{DOSCAR}和\textrm{PROCAR}文件。\textrm{PROCAR}文件包括\textrm{Kohn-Sham}轨道指定角动量$l$和磁量子数$m$的分波\textrm{DOS}，投影用的投影函数也写入了\textrm{PROCAR}文件。

执行\textrm{Bash}脚本完成计算

采用\textrm{Py4vasp}模块绘制\ch{Ni}表面的每一层的\textrm{DOS}图。%执行\textrm{Python}脚本，会给出\textrm{Py4vasp}的\textrm{DOS}作图函数，修改代码可以绘制指定层的\textrm{DOS}。相比于接近体相(第三层)的\textrm{DOS}，投影\textrm{DOS}在靠近表面部分(第五层)，会表现出更显著的交换裂分。

在\textrm{OUTCAR}文件中，还分别给出相关元素的局域电荷和磁矩信息。

不难看出，表面的第一层和第五层，电荷密度降低了，但磁性明显提升。

采用\textrm{Py4vasp}模块，可以可视化在位的自旋磁矩，有在位自旋磁矩贡献的原子位置上有箭头标注。

如果修改\textrm{INCAR}中的计算控制参数\texttt{LORBIT}=1且\texttt{RWIGS}=1.286，并再次执行\textrm{VASP}计算，注意\textrm{DOS}、局域电荷密度和磁矩的变化情况。

\subsubsection{\ch{Ni}的(100)表面的能带}
沿着高对称性方向$\Gamma-X-M-\Gamma$绘制\ch{Ni}的$(100)$表面的能带结构。

在构建表面的时候，晶体结构沿着特定方向拉升，并填充了真空层。该操作同样会显著影响倒空间的\textrm{Brillouin zone}的形状。%在下图中，可以看到\ch{Ni}的体相结构的第一\textrm{Brillouin}区，也称为\textrm{FCC}结构的\textrm{Weiger-Seitz}元胞。采用\textrm{SeeK-path}工具软件，上传体相\textrm{POSCAR}文件，可以得到高对称性$\vec k$点路径。

%采用同样的工具，可以得到表面计算所需要的\textrm{Weiger-Seitz}元胞。该\textrm{Weiger-Seitz}元胞有一个方向被压扁，而高对称性点的标记也都发生了改变。因为\textrm{Bravais}格子的格矢发生了改变，所以在选取高对称性点和方向时要比较注意。上图中，上面一张图给出表面计算所需要的\textrm{Wigner-Seitz}元胞的高对称性点。

\ch{Ni}的$(100)$面结构文件\textrm{POSCAR}

\textrm{POSCAR}文件

$\vec k$空间高对称性点及路径的选择

执行\textrm{Bash}脚本，采用前面计算完成的\textrm{CHGCAR}，完成\textrm{VASP}计算

%静态计算
%计算过程电荷密度保持不变
%自旋极化计算

采用\textrm{Py4vasp}绘制能带图，并将指定角动量量子数$l$的能带标出。

\subsubsection{\ch{Ni}的(111)表面的弛豫}
采用非自旋极化的\textrm{RMM-DIIS}方法弛豫\ch{Ni}的$(111)$表面五层结构，并确定表面能。

%在每个电子步迭代只能，在对包括电荷密度在内的物理量进行计算的时候，都需要对第一\textrm{Brillouin}积分。对于金属和导体体系，由于\textrm{Fermi}面附近能带的部分占据(称为导带)，因此$\vec k$空间积分会有不连续的函数。在$T=0\mathrm{K}$，这些位置就是阶梯函数\textrm{(Step function)}。换言之，对于导带电子的积分需要非常密集的$\vec k$空间网格剖分，这样严重消耗计算资源。计算控制参数\texttt{ISMEAR}为用户提供了不同的解决方案。其中一个近似方案就是采用\textrm{Methfessel-Paxton}方法积分，该方法基于对台阶函数的平滑近似，采用多项式逼近台阶函数的方法，得到较为精确的积分结果。多项式的截断阶数由\texttt{ISMEAR}$>N$确定。

\ch{Ni}的$(111)$表面结构文件\textrm{POSCAR}

体相\ch{Ni}的结构文件\textrm{bulk/POSCAR}

\textrm{INCAR}

\textrm{bulk/INCAR}

$\vec k$空间剖分网格文件\textrm{KPOINTS}

\textrm{bulk/KPOINTS}

在\textrm{POSCAR}中，与$(100)$面类似，体相最小重复单元含有一个原子。划分的e$(111)$面的真空层厚度为
\begin{displaymath}
	(1-0.4)\times5.774\times3.53\AA\approx12.23\AA
\end{displaymath}
因此真空层为$12.12\AA$。\texttt{selective dynamics}选项指定了两层弛豫，另外三层固定。

\textrm{INCAR}文件除了关闭自旋，与$(100)$计算完全相同，另外为了加速收敛，计算控制参数\texttt{NBANDS}确定的能带数目增加了一些。

执行\textrm{Bash}脚本

采用\textrm{Py4vasp}模块完成晶胞可视化并计算表面面积。
%边长为$a=2.5\AA$，有一个角为$120^{\circle}$的平行四边形面积为
%\begin{displaymath}
%(2.5\AA)^2\cdot\sim(120/360\times\pi)=5.41236587737\AA	
% \end{displaymath}
%或者取两个全等三角形$A=(a^2\sqrt3)/4$。\ch{Ni}的$(111)$Mi按比$(100)$面的原子堆积更紧实。

不难看出，$(111)$表面的受力更小一些，因为\ch{Ni}的$(111)$面更接近体相结构。按照同样的表面能计算方法，有非弛豫表面能$\sigma^{\mathrm{unrel}}$:
\begin{displaymath}
	\sigma^{\mathrm{unrel}}=\dfrac1{2A}(E_{\mathrm{slab}}^{\mathrm{unrel}}-NE_{\mathrm{bulk}})=122.7\mathrm{meV/AA^2}
\end{displaymath}
上面两层弛豫后的平均表面能$\langle\sigma\rangle$
\begin{displaymath}
	\langle\sigma\rangle=\dfrac1{2A}(E_{\mathrm{slab}}-NE_{\mathrm{bulk}})=122.1\mathrm{meV/AA^2}
\end{displaymath}
可得完全弛豫表面能$\sigma^{\mathrm{rel}}$为
\begin{displaymath}
	\sigma^{\mathrm{rel}}=2\langle\sigma\rangle-\sigma^{\mathrm{unrel}}=121.6\mathrm{meV/AA^2}
\end{displaymath}

在$(111)$面，弛豫引起每个面上的能量增益比$(100)$面更小，因为$(111)$面的原子更接近密堆积，所以原子间的间距更短。在\ch{Ni}的\textrm{FCC}结构中，$(111)$面的表面能最小。即使计算中还忽略了磁性的贡献，对比文献中表面能结果$\sim120\mathrm{meV/AA}$，误差也是最小的。\upcite{SD3-160080_2016}

实际计算流程中要注意
\begin{enumerate}
	\item 开始计算时要注意$\vec k$空间网格剖分(\textrm{KPOINTS}文件)的收敛测试，截断能(\texttt{ENCUT})的收敛测试；同样，真空层厚度和\textrm{Slab}厚度也应该进行收敛测试。
	\item 在\textrm{OUTCAR}文件中，会存在很大的\textrm{Pulay}应力(\textrm{external pressure}的数值)，可以通过增大截断能(\texttt{ENCUT})避免。为了对比不同计算的总能，必须设置相同的\texttt{ENCUT}的值。所有的表面能都应该在足够大的截断能，确保没有\textrm{Pulay}应力可以忽略的条件下进行比较。
	\item 再者，这里体相材料默认为给定的。当晶格常数随着交换-相关泛函发生较大变化时，应该对体相材料先进行合理的结构弛豫。
\end{enumerate}

\subsection{\ch{CO}分子在\ch{Ni}表面的吸附}
\subsubsection{\ch{CO}分子在\ch{Ni}的$(111)$面上的吸附}
完成\ch{CO}分子在五层的\ch{Ni}的$(111)$面(三层固定)上，通过\textrm{RMM-DISS}方法弛豫，并与单个分子对比结构的变化。

分子在表面吸附上是表面科学中普通却又复杂的问题。即使对于简单分子，构型空间也很大，因此高维的势能面也可以很平坦。一般地，分子可以吸附在不同方向，由于与表面的相互作用和电荷转移，原本等价的键角也会发生变化。\ch{CO}吸附是催化反应中常用的试剂，因为\ch{CO}中毒后会变成\ch{CO2}。\upcite{NJC47-20222_2023}。

吸附会显著地改变\textrm{Slab}上的电荷分布，因此系统会出现偶极矩。对周期性边界条件，有限偶极会通过真空区在势能中形成梯度。但是计算的目的是模拟单个孤立的\textrm{Slab}-吸附体系，该体系与周期坏境不发生相互作用。为了消除不希望的电场影响，有必要引入偶极校正\textrm{(dipole correction)}。该校正可以修正势能、能量和受力，并且推荐在任何具有反演对称性破缺(即允许存在偶极矩)的体系。

最困难的是准备初始结构文件。采用\textrm{ASE}模块，其中包括方便的结构生成模块\textrm{builder}，执行以下代码产生孤立的\ch{CO}分子的\textrm{POSCAR}和吸附-\textrm{Slab}-吸附物的\textrm{POSCAR}。

在表面上有不同的吸附位点:
\begin{itemize}
	\item \textcolor{red}{红点位置}:~吸附物位于顶点\textrm{(Top)}位。这是在高覆盖率\textrm{(high coverage)}\footnote{高覆盖率是表面完全被\ch{CO}覆盖的情况。}情况下最稳定的吸附位。
	\item 
\end{itemize}
w为了模拟低覆盖的情形，采用超晶胞。关于其他高对称的吸附位点，可以在\textrm{ASE}中的\textrm{suface}模块的说明。

同样地，在\textrm{POSCAR}文件中固定三层\ch{Ni}原子，允许两层表面原子弛豫。
